\renewcommand{\arraystretch}{1.4}
\begin{longtable}{c p{14cm}}
        $\R$                       & : Himpunan bilangan real                                                                 \\
        $\R_{\max}$                & : Himpunan bilangan real yang diperluas dengan elemen $\varepsilon=-\infty$              \\
        $\varepsilon$              & : Elemen nol pada aljabar max-plus, yaitu $-\infty$                                      \\
        $\oplus$                   & : Operasi penjumlahan dalam aljabar max-plus, yaitu operasi maksimum                     \\
        $\otimes$                  & : Operasi perkalian dalam aljabar max-plus, yaitu operasi penjumlahan biasa              \\
        $\bigoplus$                & : Penjumlahan berulang dalam aljabar max-plus                                            \\
        $\overline{n}$             & : Himpunan simbol $\{1,2,\ldots,n\}$                                                     \\
        $L_S(n)$                   & : Himpunan semua Latin square berordo $n$ dengan himpunan simbol $N$                     \\
        $I_n$                      & : Matriks identitas berordo $n$                                                          \\
        $P_\sigma$                 & : Matriks permutasi yang bersesuaian dengan permutasi $\sigma$                           \\
        $P_{\sigma_i^A}$           & : Matriks permutasi untuk posisi simbol ke-$i$ pada $A$                                  \\
        $P_{\sigma_j^B}$           & : Matriks permutasi untuk posisi simbol ke-$j$ pada $B$                                  \\
        $P^{(t)}$                  & : Matriks permutasi sirkulan-$t$ atas                                                    \\
        $P_{(t)}$                  & : Matriks permutasi sirkulan-$t$ bawah                                                   \\
        $C^{(t)}$                  & : Matriks sirkulan-$t$ atas                                                              \\
        $C_{(t)}$                  & : Matriks sirkulan-$t$ bawah                                                             \\
        $\sigma_i^A$               & : Permutasi yang menyatakan posisi simbol ke-$i$ pada $A$                                \\
        $\sigma_j^B$               & : Permutasi yang menyatakan posisi simbol ke-$j$ pada $B$                                \\
        $C_{S_n}(\sigma)$          & : Centralizer dari permutasi $\sigma$ di $S_n$                                           \\
        $C_{S_n}(H)$               & : Centralizer dari subgrup $H$ di $S_n$                                                  \\
        $H_A$                      & : Subgrup yang dibangkitkan oleh permutasi-permutasi penyusun $A$                        \\
        $\Sigma_A$                 & : Himpunan permutasi penyusun $A$, yaitu $\{\sigma_1^A,\sigma_2^A,\ldots,\sigma_n^A\}$   \\
        $B^*$                      & : Matriks kandidat parsial untuk membentuk kandidat Latin square $B$                     \\
        $b^*_{rs}$                 & : Entri baris ke-$r$ dan kolom ke-$s$ dari matriks kandidat parsial $B^*$                \\
        $\mathcal{P}(A)$           & : Himpunan kandidat Latin square yang komutatif dengan $A$ yang diperoleh dari algoritma \\
        $\Gamma(\sigma)$           & : Representasi koordinat dari permutasi $\sigma$, yaitu $\{(r,\sigma(r))\mid r\in N\}$   \\
        $\mathcal{U}_{\ge v}^{AB}$ & : Himpunan superlevel dari $A\otimes B$ pada level $v$                                   \\
        $\mathcal{U}_{\ge v}^{BA}$ & : Himpunan superlevel dari $B\otimes A$ pada level $v$                                   \\
        $C_n$                      & : Grup siklik berordo $n$                                                                \\
        $V_4$                      & : Grup Klein, yaitu grup yang isomorfik dengan $C_2\times C_2$                           \\
\end{longtable}